\documentclass[]{article}

\usepackage{amsmath}
\usepackage{multirow}
\usepackage{natbib}
\usepackage{graphicx} 


\begin{document}

% Auto-generated figures file
% \usepackage{graphicx}

\begin{figure}[h!]
    \centering
    \includegraphics[width=0.5\textwidth]{../input_files/plots/step_1_energy_correction_summary_1775399991.png}
    \caption{Diagnostic plots comparing the raw measured energy (blue) and the corrected energy (red) for Oscillators 1, 10, and 20. The raw energy signal exhibits a persistent plateau at late times due to a noise-induced bias. This noise floor is effectively eliminated by the deconvolution framework, resulting in corrected energy trajectories that show a clean exponential decay toward zero, consistent with the theoretical model of a damped system.}
    \label{fig:image0}
\end{figure}

\begin{figure}[h!]
    \centering
    \includegraphics[width=0.5\textwidth]{../input_files/plots/step_2_validation_plots_1775400009.png}
    \caption{Validation of the noise deconvolution framework for 20 oscillators. (Left) The observed damping rate ($\gamma_{obs}$), derived from noise-corrected energy data, is plotted against the theoretical damping rate ($\gamma_{theory}$). The data points show a strong linear correlation, clustering around the ideal identity line (dashed). (Right) A histogram of the residuals ($\gamma_{obs} - \gamma_{theory}$) shows a distribution centered near zero, indicating that the deconvolution method successfully removes systematic bias in the damping rate estimation.}
    \label{fig:image1}
\end{figure}



\end{document}
                